% !TEX program = xelatex
\documentclass[12pt]{article}

\usepackage{amsmath}
\usepackage{setspace}
\usepackage[unicode]{hyperref}
\usepackage{graphicx}
\usepackage{subcaption}
\usepackage[labelfont={bf,large},font={large}]{caption}
\usepackage{xepersian}
\settextfont{XB Niloofar}



\title{گزارش پیشرفت}

% \date{\vspace{-8ex}\today}

\date{\vspace{-6ex}\normalsize بهمن 1402}

\setstretch{1.1}


\begin{document}
% title

\vspace{-10ex}
\centerline{به نام او}
{\let\newpage\relax\maketitle}
\textbf{
    \large
    \centerline{عنوان پیشنهاده: ناهنجاری جهتی در داده‌های بزرگ مقیاس کیهانی}
}
% end title

\thispagestyle{empty}



در نیم سال گذشته تمرکز ما بر روی بازتولید ناهنجاری‌های گزارش شده در تابش زمینه کیهانی بود.
یکی از این ناهنجاری‌ها، ناهنجاری دوقطبی در نقشه واریانس دمایی است.
\lr{[Planck VII 2018]}
روش به کار رفته در پژوهش فوق،
برای اولین بار در
\lr{[Akrami et.al 2014]}
انجام شد.
ما طی جلساتی با آقای دکتر اکرمی مقاله 
\lr{[Akrami et.al 2014]}
را به طور کامل بازتولید کردیم و نتایج آن هم که در سازگاری با نتایج پیشین بود در ادامه آمده است.
این روش به شرح زیر است:
\subsection*{
    ناهنجاری واریانس موضعی
}

ما 3072 دیسک
(عرقچین)
با اندازه‌های مختلف به مرکز پیکسل‌های یک نقشه‌ی 
\lr{HEALPix}
 با 
$N_{side} = 16$
را در نظر می‌گیریم
(این تعداد، همان تعداد مقاله یاد شده است).

برای هر عرقچین و نقشه‌ی آسمانی، واریانس دمایی را تنها برای پیکسل‌هایی که ماسک نشده‌اند،
محاسبه می‌کنیم؛ هر عرقچینی که بیش از 90 درصد آن ماسک شده باشد فاقد اعتبار دانسته، در محدوده‌ی ماسک در نظر گرفته می‌شود.
با نسبت دادن این مقادیر واریانس به پیکسل‌های نقشه‌ی 
$N_{side} = 16$
یک نقشه‌ی واریانس موضعی
(که دارای ماسک هم هست)
به دست می‌آید.
این کار را هم برای خود داده 
\lr{CMB}
و هم برای 1000 شبیه‌سازی که با الگوریتم یکسان تهیه شده‌اند
(و نویز هم به آنها اضافه شده است)
انجام می‌دهیم و به این ترتیب
1000 + 1
نقشه‌ی واریانس موضعی خواهیم داشت.

میانگین و واریانس این نقشه‌های واریانس، خود مهم هستند چرا که باید بررسی کرد که مقادیر 
به دست آمده از 
\lr{CMB}
یا هر یک از شبیه‌سازی‌ها به چه میزان از مقدار میانگین فاصله دارد و چقدر نامحتمل است.
برای به دست آوردن میانگین و واریانس انتظاری از هر عرقچین، 
از آنجا که نویز هم به نقشه‌ها اضافه شده است، مکان عرقچین‌ها اهمیت پیدا می‌کنند
چرا که با افزودن نویز می‌توان آن دسته از عرقچین‌هایی که مکان یکسانی دارند
(در همه 1000+1 نقشه)،
را از یک آنسامبل دانست.
پس با محاسبه‌ی میانگین و واریانس عرقچین‌های هم‌مکان
میانگین و واریانس آنسامبلی نقشه‌ی واریانس را به دست می‌آوریم.
حال میانگین را از همه عرقچین‌های هم‌مکان کم می‌کنیم.

اما برای مقایسه هنوز یک گام دیگر باید برداریم.
عرقچین‌هایی که مکان‌های مختلفی دارند باید وزن‌های متفاوتی داشته باشند
چرا که میزان پیکسل‌های ماسک شده در آنها متفاوت است.
حال هر قدر تعداد پیکسل‌های ماسک شده بیشتر باشد، افت و خیز در یک عرقچین بیشتر است.
عکس واریانس آنسابلی را به عنوان وزن گرفته و در مقدار بالا ضرب می‌کنیم:

\begin{align}
    \nonumber
    \frac{1}{Var(Var)} \big(Var(T) - <Var(T)>\big)
\end{align}

حال با در دست داشتن این مقدار در تمام پیکسل‌ها، 
یک دوقطبی به نقشه برازش می‌کنیم که جهت و دامنه ناهنجاری را به ما می‌دهد.
مقایسه دامنه دوقطبی 
\lr{CMB}
با شبیه‌سازی‌ها برای اندازه‌های مختلف عرقچین
(دیسک)
به ما نشان می‌دهد که
\lr{CMB}
به چه میزان بهنجار یا ناهنجار است. در
\hyperref[plt_01]{نمودار 1}
\lr{p-value}
به دست آمده در مقاله و آنچه ما به دست آورده‌ایم، آمده است:

همچنین اگر نقشه‌ی واریانس 
(که مقدار میانگین از آن کم شده و وزن دهی شده است)
را بر حسب هماهنگ‌های کروی بسط بدهیم و از ضرایب آن 
$C_l$
(یا همان واریانس ضرایب 
$a_{lm}$
)
را محاسبه کنیم، خواهیم دید که جمله مربوط به 
$l=1$
از دیگر مقیاس‌ها بزرگتر است
(
\hyperref[plt_02]{نمودار 2}
)

\subsection*{ایده پیش رو:}

از آنجا که سنجه 
$C_l$
تنها وابسته به جدایی زاویه‌ای است و جهت در آن بی معنی است،
ما تلاش داریم این کار را به گونه‌ای دیگر تکرار کنیم که در آن جهت‌های مختلف تمییز داده شوند.
برای این منظور ما از هندسه‌ی نواری در مکان‌ها و جهات مختلف بهره می‌بریم.
نوارها به ما این امکان را می‌دهند که بتوانیم ضرایب بسط 
$\sum a_{lm} Y_{lm}$
را برای 
$l=0$
در همه جهات مقایسه کنیم و بررسی کنیم که آیا جهت مرجحی برای اینگونه ناهنجاری وجود دارد یا خیر.
\\
\\
\\

\noindent
با احترام

\noindent
محمدحسین جمشیدی

\newpage

\begin{thebibliography}{9}
    \begin{latin}
        
        \bibitem{akrami2014power}
        Akrami, Y and Fantaye, Y and Shafieloo, A and Eriksen, HK and Hansen, Frode Kristian and Banday, Anthony J and G{\'o}rski, Krzysztof M,
        \textit{
            Power asymmetry in WMAP and Planck temperature sky maps as measured by a local variance estimator
            }
            ,
            The Astrophysical Journal Letters,
            2014.
        \bibitem{akrami2020planck},
        Akrami, Yashar and Ashdown, M and Aumont, Jonathan and Baccigalupi, Carlo and Ballardini, M and Banday, Anthony J and Barreiro, RB and Bartolo, Nicola and Basak, S and Benabed, K and others,
        \textit{Planck 2018 results-VII. Isotropy and statistics of the CMB},
        Astronomy \& Astrophysics,
        2020.
    
    \end{latin}
\end{thebibliography}

\begin{figure}
    \begin{center}
        \includegraphics[height=8cm]{DipoleAmplitude.png}
        \includegraphics[height=6cm]{DipoleAmplitude_akrami.PNG}
        \caption[]
        {
            \normalsize
            مقایسه دامنه‌های دوقطبی بازتولید شده 
            (بالا)
            و  
            \lr{Akrami et.al 2014} (پایین)

            اندک تفاوت میان نتایج به علت استفاده از سری داده‌های متفاوت 
            پلانک است که در سال‌های مختلف منتشر شده
            }
    \end{center}
    \label{plt_01}
\end{figure}

\begin{figure}
    \begin{center}
        \includegraphics[height=8cm]{C_l.png}
        \includegraphics[height=6cm]{C_l_akrami.PNG}
        \caption[]
        {
            \normalsize
            مقایسه $C_l$ بازتولید شده 
            (بالا)
            و  
            \lr{Akrami et.al 2014} (پایین)
            }
    \end{center}
    \label{plt_02}
\end{figure}

\end{document}